\documentclass[aspectratio=169, xcolor=dvipsnames]{beamer}

\AtBeginSection[]
{
    \begin{frame}
        \frametitle{Visão Geral}
        \tableofcontents[currentsection]
    \end{frame}
}

\usefonttheme{professionalfonts} % using non standard fonts for beamer
\usefonttheme{serif} % default family is serif
\usepackage{fontspec}
\setmainfont{XCharter}[
    UprightFont    = *-Roman,
    BoldFont       = *-Bold,
    ItalicFont     = *-Italic,
    BoldItalicFont = *-BoldItalic,
]

\usepackage{hyperref}
\usepackage{graphicx} % Allows including images
\usepackage{booktabs} % Allows the use of \toprule, \midrule and \bottomrule in tables
\input{./slides/SimplePlusTheme/beamerthemeSimplePlus.sty}
\input{./slides/SimplePlusTheme/beamerinnerthemeSimplePlus.sty}
\input{./slides/SimplePlusTheme/beamerfontthemeSimplePlus.sty}
\input{./slides/SimplePlusTheme/beamercolorthemeSimplePlus.sty}

\usepackage{listings}
\setbeamercovered{transparent}

\usepackage{biblatex}
\addbibresource{slides/reference.bib}
\renewcommand*{\bibfont}{\footnotesize}
% Tamanho dos slides de "Leituras Recomendadas": maior que a bibliografia
% completa (\bibfont acima, que rende \tiny via a redefinição de \footnotesize)
% para destacar as leituras curadas.
\newcommand{\leiturasize}{\small}

\usepackage[brazilian ]{babel}
\usepackage{csquotes}
\usepackage{pgfplots}
\pgfplotsset{compat=1.18}

\usepackage{emoji}
\usepackage{amsmath}
\usepackage[xcharter,bigdelims,vvarbb]{newtxmath}
% newtxmath's operators font requests the fontspec family `XCharter(0)' in OT1;
% without these substitutions math digits and operator names silently fall back
% to Computer Modern (LaTeX Font Warning: Font shape `OT1/XCharter(0)/m/n' undefined).
\DeclareFontFamily{OT1}{XCharter(0)}{}
\DeclareFontShape{OT1}{XCharter(0)}{m}{n}{<->ssub * XCharter-TLF/m/n}{}
\DeclareFontShape{OT1}{XCharter(0)}{m}{it}{<->ssub * XCharter-TLF/m/it}{}
\DeclareFontShape{OT1}{XCharter(0)}{b}{n}{<->ssub * XCharter-TLF/b/n}{}
\DeclareFontShape{OT1}{XCharter(0)}{bx}{n}{<->ssub * XCharter-TLF/b/n}{}

\usepackage{bm}
\usepackage[table]{xcolor}
\usepackage{xcolor}
\usepackage{algpseudocode}

\usepackage{cancel}
\usepackage{ulem}

\usetikzlibrary{shapes}
\usetikzlibrary{arrows}
\usetikzlibrary {arrows.meta}
\usetikzlibrary{calc,positioning}
\usepgfplotslibrary{fillbetween}
\usetikzlibrary{angles,quotes}
\usetikzlibrary{tikzmark}


\definecolor{PastelRed}{HTML}{FFC1CC}
\definecolor{PastelBlue}{HTML}{C2F0FF}
\definecolor{PastelBlue2}{HTML}{3182BD}

\definecolor{PastelPink}{HTML}{FFCBE1}
\definecolor{PastelGreen}{HTML}{D6E5BD}
\definecolor{PastelYellow}{HTML}{F9E1A8}
\definecolor{PastelBlue}{HTML}{BCD8EC}
\definecolor{PastelPurple}{HTML}{DCCCEC}
\definecolor{PastelOrange}{HTML}{FFDAB4}


\renewcommand{\footnotesize}{\tiny}

\newcommand{\mespor}{
  \ifcase\month
    \or janeiro
    \or fevereiro
    \or março
    \or abril
    \or maio
    \or junho
    \or julho
    \or agosto
    \or setembro
    \or outubro
    \or novembro
    \or dezembro
  \fi
}
% \usepackage{csquotes}
% \usepackage{minted}
% \usemintedstyle{xcode}
\definecolor{vscLightBackground}{HTML}{FFFFFF}
\definecolor{vscLightText}{HTML}{000000}
\definecolor{vscLightGreen}{HTML}{008000}        % For Comments AND Numbers
\definecolor{vscLightRed}{HTML}{A31515}          % For Strings
\definecolor{vscLightTeal}{HTML}{267F99}
\definecolor{vscBlue}{HTML}{0000FF}% For main keywords (if, for, def...)
\definecolor{vscLightPytorch}{HTML}{800080}      % For PyTorch keywords (dark magenta)

\lstdefinestyle{vscode-light}{
    language=Python,
    backgroundcolor=\color{vscLightBackground},
    basicstyle=\ttfamily\scriptsize\color{vscLightText},
    keywordstyle=\color{vscLightTeal},
    commentstyle=\color{vscLightGreen},
    stringstyle=\color{vscLightRed},
    % Style for PyTorch keywords (using keyword group 2)
    keywordstyle=[2]\color{vscLightPytorch}, 
    morekeywords={self, True, False, None},
    morekeywords=[2]{
        torch, Tensor, device, dtype, to, nn, optim, autograd, utils, data,
        Module, Linear, Conv2d, ReLU, CrossEntropyLoss, Adam, SGD,
        Dataset, DataLoader, randn, zeros, ones, arange, cat, stack,
        view, reshape, sum, mean, max, min, matmul, no_grad, backward
    },
    keywordstyle=[3]\color{vscBlue}, 
    morekeywords=[3]{
        model, loss, y, y_hat, x , optimizer
    },
    frame=single,
    tabsize=1,
    showstringspaces=false,
    breaklines=true,
    captionpos=b,
    extendedchars=true,
    numbers=left,
}

\usepackage[cache=true]{minted}
% minted v3+ (TeX Live 2024+) auto-detects the output directory.
% minted v2 (Ubuntu apt TeX Live 2023) needs it set explicitly to
% match latexmkrc's $out_dir = 'build'.
\makeatletter
\@ifpackagelater{minted}{2024/01/01}{}{%
  \def\minted@outputdir{build/}%
}
\makeatother
\setminted{
    fontsize=\scriptsize,
    style=vs,
    breaklines=true,
    numbers=left,
    numbersep=5pt,
    frame=single,
}
\providecommand{\citeauthoryear}[1]{(\citeauthor{#1},~\citeyear{#1})}

\providecommand{\thetav}{\bm{\theta}}

% vetor coluna 2D com colchetes
\providecommand{\cvec}[2]{\begin{bmatrix} #1 \\ #2 \end{bmatrix}}

\providecommand{\varvector}[1]{\mathbf{#1}}
% instance
\providecommand{\X}{\varvector{X}}
\providecommand{\mlinstance}{\varvector{X}}
\providecommand{\mlinstancei}{\varvector{X}^{(i)}}
\providecommand{\mlinstancex}[1]{\varvector{X}^{(#1)}}

% target
\providecommand{\y}{\varvector{y}}
\providecommand{\mltarget}{\varvector{y}}
\providecommand{\mltargeti}{\varvector{y}^{(i)}}
\providecommand{\mltargetx}[1]{\varvector{y}^{(#1)}}

%hats
\providecommand{\yhat}{\hat{\varvector{y}}}
\providecommand{\mltargethat}{\hat{\varvector{y}}}
\providecommand{\mltargethati}{\hat{\varvector{y}}^{(i)}}
\providecommand{\mltargethatx}[1]{\hat{\varvector{y}}^{(#1)}}

% linear reg
\providecommand{\mllinearregression}{\mltargethat=\mlinstance\bm{\theta}}

\providecommand{\tikzarrowcolor}[3]{%
  \begin{tikzpicture}[remember picture, overlay]
    % Arrow body
    \path[fill=#3] ([shift={(#1,#2)}]current page.south west) 
      rectangle ++(0.3, -0.2);
    % Arrow head
    \path[fill=#3] ([shift={(#1+0.3,#2+0.1)}]current page.south west) -- 
         ++(0, -0.4) -- 
         ++(0.2, 0.2) -- 
         cycle;
  \end{tikzpicture}%
}

\providecommand{\tikzarrow}[2]{%
    \tikzarrowcolor{#1}{#2}{Red}
}



\DeclareMathOperator*{\argmax}{argmax}
\DeclareMathOperator*{\argmin}{argmin}



%----------------------------------------------------------------------------------------
%    TITLE PAGE
%----------------------------------------------------------------------------------------

\title{Deep Learning}
\subtitle{\textit{Backpropagation}}

\author{Lucas Silveira Kupssinskü}

\institute
{ Machine Learning Theory and Applications Laboratory \\ PUCRS % Your institution for the title page
}
\date{\mespor de \the\year}
% Date, can be changed to a custom date

%----------------------------------------------------------------------------------------
%    PRESENTATION SLIDES
%----------------------------------------------------------------------------------------

\begin{document}
  \begin{frame}
    % Print the title page as the first slide
    \titlepage
  \end{frame}

  \begin{frame}{Overview}
    % Throughout your presentation, if you choose to use \section{} and \subsection{} commands, these will automatically be printed on this slide as an overview of your presentation
    \tableofcontents
  \end{frame}
  \section{Introdução}

  \begin{frame}{MLP}
    \begin{columns}
      \begin{column}{.65\linewidth}
        \begin{itemize}
          \item Um MLP é um aproximador universal de funções\footnotemark\textsuperscript{,}\footnotemark

          \item Contudo, o PLA (\textit{Perceptron Learning Algorithm}) não funciona
            para um MLP
        \end{itemize}
        \centering
        \resizebox{.45\linewidth}{!}{
        \begin{tikzpicture}
          \pgfmathsetmacro{\ex}{0.5} \pgfmathsetmacro{\ey}{1.} \pgfmathsetmacro{\eradius}{0.5}
          \pgfmathsetmacro{\yone}{1.5} \pgfmathsetmacro{\ytwo}{0} \pgfmathsetmacro{\ythree}{-1.5}
          % \draw[help lines] (-2,-2) grid (3,3);
          % first
          \fill (-2, \yone) circle[x radius=0.15, y radius=0.15];
          \node at (-2, \yone+0.5) {$x_{0}=1$};

          % second
          \fill (-2, \ytwo) circle[x radius=0.15, y radius=0.15];
          \node at (-2, \ytwo+0.5) {$x_{1}$};

          % third
          \fill (-2, \ythree) circle[x radius=0.15, y radius=0.15];
          \node at (-2, \ythree+0.5) {$x_{2}$};

          \draw[thick] (-2, \yone) -- (\ex, \ey-1);
          \draw[thick] (-2, \ytwo) -- (\ex, \ey-1);
          \draw[thick] (-2, \ythree) -- (\ex, \ey-1);
          \node at (-1.5,{\yone}) {\small \textcolor{blue}{$-1.5$}};
          \node at (-1.5,{\ytwo+0.2}) {\small \textcolor{blue}{$1$}};
          \node at (-1.5,{\ythree+0.6}) {\small \textcolor{blue}{$-1$}};
          \draw[thick] (\ex, 0) -- (\ex+2, 0);
          \filldraw[color=black, fill=PastelBlue] (\ex,\ey-1) circle (0.5);
          \node[font=\tiny] at (\ex, 0) {$\varphi(z^{(i)})$};

          \node at (-1.5,{\yone-1}) {\small \textcolor{green}{$-1.5$}};
          \node at (-1.5,{\ytwo-0.2}) {\small \textcolor{green}{$-1$}};
          \node at (-1.5,{\ythree-0.3}) {\small \textcolor{green}{$1$}};
          \draw[thick] (-2, \yone) -- (\ex, \ythree);
          \draw[thick] (-2, \ytwo) -- (\ex, \ythree);
          \draw[thick] (-2, \ythree) -- (\ex, \ythree);

          \draw[thick] (\ex,\ythree) -- (\ex+2, 0);

          \filldraw[color=black, fill=PastelGreen] (\ex,\ythree) circle (0.5);

          \node[font=\tiny] at (\ex, \ythree) {$\varphi(z^{(i)})$};%
          \draw[thick] (-2, \yone) -- (\ex+2, 0);
          \node at (1.7,{\yone-1}) {\small \textcolor{red}{$0.5$}};
          \node at (1.7,{\ytwo-0.2}) {\small \textcolor{red}{$1$}};
          \node at (1.7,{\ytwo-0.8}) {\small \textcolor{red}{$1$}};

          \draw[thick, ->] (\ex+2, 0) -- (\ex+3, 0);
          \filldraw[color=black, fill=PastelRed] ({\ex+2},{\ey-1}) circle (0.5);
          \node[font=\tiny] at (\ex+2, 0) {$\varphi(z^{(i)})$};%
        \end{tikzpicture}
        }
      \end{column}
      \begin{column}{.35\linewidth}
        \centering
        \resizebox{\linewidth}{!}{
        \begin{tikzpicture}
          \begin{axis}[
            legend pos=south east, % <-- this line controls the width
            xlabel={$x_{1}$},
            ylabel={$x_{2}$},
            xmin=-1.5,
            xmax=1.5,
            ymin=-1.5,
            ymax=1.5,
            xtick={-1, 0, 1},
            ytick={-1, 0, 1},
            ymajorgrids=true,
            xmajorgrids=true,
            grid style=dashed,
          ]
            \path[name path=bottom]
              (\pgfkeysvalueof{/pgfplots/xmin},\pgfkeysvalueof{/pgfplots/ymin}) -- (\pgfkeysvalueof{/pgfplots/xmax},\pgfkeysvalueof{/pgfplots/ymin});

            \path[name path=top]
              (\pgfkeysvalueof{/pgfplots/xmin},\pgfkeysvalueof{/pgfplots/ymax}) -- (\pgfkeysvalueof{/pgfplots/xmax},\pgfkeysvalueof{/pgfplots/ymax});

            \addplot [ name path=g, domain=-1.5:1.5, samples=2, color=black,
            dashed, thick, ] {x+1}; \addplot[PastelBlue, opacity=0.3] fill between
            [of=g and top]; \addplot[PastelRed, opacity=0.3] fill between [of=g and
            bottom]; \addplot [ name path=f, domain=-1.5:1.5, samples=2, color=black,
            dashed, thick, ] {x-1};

            \addplot[PastelRed, opacity=0.3] fill between [of=f and top]; \addplot[PastelBlue,
            opacity=0.3] fill between [of=f and bottom]; \addplot [ name path=f,
            domain=-1.5:1.5, samples=2, color=black, dashed, thick, ] {x-1};
            \addplot [ name path=g, domain=-1.5:1.5, samples=2, color=black, dashed,
            thick, ] {x+1}; \addplot[PastelBlue, opacity=0.3] fill between [of=g
            and top]; \addplot[PastelRed, opacity=0.3] fill between [of=f and g];
            \addplot[PastelBlue, opacity=0.3] fill between [of=f and bottom]; \addplot[
            only marks, ultra thick, color=red, mark=-, mark options={scale=3} ]
            coordinates { (-1,-1)(1,1) }; \addplot[ only marks, ultra thick,
            color=blue, mark=+, mark options={scale=3} ] coordinates
            { (-1,1)(1,-1) };
          \end{axis}
        \end{tikzpicture}}
      \end{column}
    \end{columns}
    \footnotetext[1]{\fullcite{hornik1989multilayer}} \footnotetext[2]{\fullcite{cybenko1989approximationsigmoid}}
  \end{frame}

  \begin{frame}{Inverno das Redes Neurais}
    \begin{itemize}
      \item Com a publicação de Minsky e Papert (1969) enfatizando as limitações
        do Perceptron\cite{minsky1969perceptron} o interesse em redes neurais
        diminuiu.
        \begin{itemize}
          \item Pouca verba de pesquisa

          \item Por que pesquisar um modelo linear enquanto temos modelos lineares
            melhores (Regressão Linear) e diversos modelos não-lineares?

          \item MLP resolve problemas não-lineares, mas ninguém sabe como
            treinar um MLP
        \end{itemize}

      \item Foi apenas 20 anos depois que começaram a surgir as primeiras publicações
        com um método de treinamento para MLPs e redes neurais em geral\footfullcite{rumelhart1986learningbackpropextended}
        \begin{itemize}
          \item Hinton não alega ter inventado o \textit{backpropagation}, mas ele
            foi um dos principais responsáveis pela sua popularização
        \end{itemize}
    \end{itemize}
  \end{frame}

  \section{Treinamento de um MLP}
  \begin{frame}
    \frametitle{Treinando um MLP}
    \begin{itemize}
      \item Vamos calcular a mão uma iteração de treino em um MLP

      \item \textbf{Ideia:} Aplicar a descida de gradiente

      \item \textbf{Função de custo:} \textit{Half Mean Squared Error}\footnote{O fator $1/2$ é convenção: simplifica a derivada do quadrático.}
        \[
          J(\thetav)=\frac{1}{2N}\sum_{i}\left(\y^{(i)}-\yhat^{(i)}\right)^{2}
        \]

      \item \textbf{Regra de atualização:}
        \[
          \thetav_{t+1}=\thetav_{t}-\eta\nabla_{\thetav} J
        \]
    \end{itemize}
    \vfill
    \textbf{Premissas para simplificar o problema:}
    \begin{itemize}
      \item Apenas 2 camadas

      \item 1 neurônio em cada camada

      \item Desconsideramos o bias

      \item Temos apenas um exemplo $\y=1, x_{1}=1$
    \end{itemize}
  \end{frame}

  \begin{frame}{Treinando um MLP}{Notação e Definições}
    \begin{itemize}
      \item $l\in [0,1,2,...,L]$ é o índice da camada do MLP ($L$ é o número de camadas)

      \item $\X$ é a entrada do MLP (também pode ser chamado de $\mathbf{a}^{(0)}$)

      \item $\y$ é a variável alvo

      \item $\theta^{(l)}_{ij}$ é o peso da camada $l$ que conecta a entrada $j$
        a saída $i$

      \item $\mathbf{z}^{(l)}$ é a \textbf{\alert{pré-ativação}} da unidade, $\mathbf{z}^{(l)}=
        \mathbf{a}^{(l-1)}\mathbf{\Theta}^{(l)^\top}$

      \item $\mathbf{a}^{(l)}$ é a \textbf{\alert{ativação}} (saída) da unidade
        $\mathbf{a}^{(l)}=\varphi(\mathbf{z}^{(l)})$
    \end{itemize}
  \end{frame}

  \begin{frame}{Treinando um MLP}{A função de Ativação}
    \begin{columns}
      \begin{column}{.4\linewidth}
        \textbf{Função Sinal:}
        \[
          \varphi(z)=
          \begin{cases}
            1  & z\geq 0 \\
            -1 & z <0
          \end{cases}
        \]
        \visible<2>{\begin{alertblock}{Problemas da função sinal}
          \begin{itemize}
            \item Não é diferenciável em $z=0$

            \item Sua derivada é 0 em todos os outros pontos
          \end{itemize}
          Isso impede o aprendizado via descida de gradiente.
        \end{alertblock}}
      \end{column}
      \begin{column}{.6\linewidth}
        \resizebox{\linewidth}{!}{
        \begin{tikzpicture}[font=\normalsize]
    \begin{axis}[
        %scale only axis,
        width=\textwidth,
        legend pos=south east, % <-- this line controls the width
        xlabel={$z$},
        ylabel={$\text{sinal}(z)$},
        xmin=-2,
        xmax=2,
        ymin=-1.5,
        ymax=1.5,
        xtick={-1.5,-1,-0.5,0,0.5,1,1.5},
        ytick={-1.5,-1,-0.5,0,0.5,1,1.5},
        ymajorgrids=true,
        xmajorgrids=true,
        grid style=dashed,
        ylabel near ticks,
        xlabel near ticks,
        ylabel style={rotate=-90},
    ]
        \addplot[ only marks, color=PastelBlue2, mark=*, ultra thick, mark
        options={scale=1.5, fill=white} ] coordinates { (0,-1) }; \addplot[ only
        marks, color=PastelBlue2, mark=*, ultra thick, mark options={scale=1.5}
        ] coordinates { (0, 1) }; \addplot[ color=PastelBlue2,
        % line width=1pt,
        ultra thick, ] coordinates { (0,1)(2.5,1) }; \addplot[ color=PastelBlue2,
        % line width=1pt,
        ultra thick, ] coordinates { (-2.5,-1)(0,-1) };
    \end{axis}
\end{tikzpicture}
        \begin{tikzpicture}[font=\normalsize]
    \begin{axis}[
        %scale only axis,
        width=\textwidth,
        legend pos=south east, % <-- this line controls the width
        xlabel={$x$},
        ylabel={$\frac{\partial \text{sinal}(z)}{\partial z}$},
        xmin=-2,
        xmax=2,
        ymin=-1.5,
        ymax=1.5,
        xtick={-1.5,-1,-0.5,0,0.5,1,1.5},
        ytick={-1.5,-1,-0.5,0,0.5,1,1.5},
        ymajorgrids=true,
        xmajorgrids=true,
        grid style=dashed,
        ylabel near ticks,
        xlabel near ticks,
        ylabel style={rotate=-90},
    ]
        \addplot [ only marks, color=PastelBlue2, mark=*, mark options={scale=1.5, fill=white},
        ultra thick, ] coordinates { (0,0) }; \addplot[ultra thick, color=PastelBlue2,
        samples=2, smooth,domain=0:6, name path=three] coordinates {(0,0)(0,2)};

        \addplot [ domain=-2:2, samples=2, color=PastelBlue2, ultra thick, ] {0};
    \end{axis}
\end{tikzpicture}

        }
      \end{column}
    \end{columns}
  \end{frame}

  \begin{frame}{Treinando um MLP}{A função de Ativação}
    \begin{columns}
      \begin{column}{.3\linewidth}
        \textbf{Função Sigmoide:}
        \[
          \varphi(z)= \sigma(z)=\frac{1}{1+e^{-z}}
        \]
        \begin{exampleblock}{Solução}
          A função sigmoide é uma versão ``suave'' da função sinal e é diferenciável
          em todo seu domínio.
        \end{exampleblock}
      \end{column}
      \begin{column}{.7\linewidth}
        \resizebox{\linewidth}{!}{
        \input{slides/img/tikz/activationfunction/sigmoidgraph.tikz}
        \begin{tikzpicture}[font=\normalsize]
    \begin{axis}[
        %scale only axis,
        width=\textwidth,
        legend pos=south east, % <-- this line controls the width
        xlabel={$z$},
        ylabel={$\frac{\partial \text{\sigma}(z)}{\partial z}$},
        xmin=-6,
        xmax=6,
        ymin=-.2,
        ymax=1.2,
        xtick={-5,-4,-3,-2,-1,0,1,2,3,4,5},
        ytick={-1, -0.75,-0.5,-0.25,0,0.25,0.5,0.75,1},
        ymajorgrids=true,
        xmajorgrids=true,
        grid style=dashed,
        ylabel near ticks,
        xlabel near ticks,
        ylabel style={rotate=-90},
    ]
        \addplot [ domain=-6:6, samples=50, color=PastelBlue2, ultra thick, ]
        {1/(1+exp(-\x))*(1-1/(1+exp(-\x)))};
    \end{axis}
\end{tikzpicture}
        }
      \end{column}
    \end{columns}
  \end{frame}

  \begin{frame}{As Etapas do Treinamento de uma Rede Neural}
    A cada iteração de treinamento, seguimos quatro passos:

    \begin{columns}
      \begin{column}{.6\linewidth}
        \begin{enumerate}
          \item \tikzmark{forward} \textbf{Forward Pass:} Computar todas as ativações
            da rede.

          \item \tikzmark{loss}\textbf{Loss Function:} Computar o valor da
            \textit{Loss Function} que quantifica o quão diferente as saídas são
            em relação à variável alvo.

          \item \tikzmark{backward}\textbf{Backward Pass:} Computar o gradiente
            da \textit{Loss Function} em relação aos parâmetros da rede.

          \item \tikzmark{optimizer} \textbf{Optimizer:} Usar os gradientes para
            atualizar os pesos.
        \end{enumerate}
      \end{column}
      \begin{column}{.4\linewidth}
        \centering
        \resizebox*{\linewidth}{!}{ \lstinputlisting[style=vscode-light,linewidth=1.05\textwidth]{slides/code/01_mintraining.py}%
        }
      \end{column}
    \end{columns}
    \begin{tikzpicture}[overlay, remember picture]
      \pgfmathsetmacro{\x}{1.1} \pgfmathsetmacro{\y}{0.15}\pgfmathsetmacro{\len}{0.5}
      \draw[->, ultra thick, red]
        ($(pic cs:forward) +(-\x, +\y)$) -- ($(pic cs:forward) +(-\x+\len,+\y)$);
      \draw[->, ultra thick, PastelBlue2]
        ($(pic cs:loss) +(-\x, +\y)$) -- ($(pic cs:loss) +(-\x+\len,+\y)$);
      \draw[->, ultra thick, PastelPink]
        ($(pic cs:backward) +(-\x, +\y)$) -- ($(pic cs:backward) +(-\x+\len,+\y)$);
      \draw[->, ultra thick, PastelGreen]
        ($(pic cs:optimizer) +(-\x, +\y)$) -- ($(pic cs:optimizer) +(-\x+\len,+\y)$);
    \end{tikzpicture}
    \begin{tikzpicture}[remember picture, overlay]
      \pgfmathsetmacro{\x}{11.65} \pgfmathsetmacro{\y}{1.8}
      \draw[->, ultra thick, red] (\x+1, \y) -- (\x, \y);
    \end{tikzpicture}%
    \begin{tikzpicture}[remember picture, overlay]
      \pgfmathsetmacro{\x}{13.} \pgfmathsetmacro{\y}{1.5}
      \draw[->, ultra thick, PastelBlue2] (\x+1, \y) -- (\x, \y);
    \end{tikzpicture}%
    \begin{tikzpicture}[remember picture, overlay]
      \pgfmathsetmacro{\x}{11.65} \pgfmathsetmacro{\y}{1.2}
      \draw[->, ultra thick, PastelPink] (\x+1, \y) -- (\x, \y);
    \end{tikzpicture}%
    \begin{tikzpicture}[remember picture, overlay]
      \pgfmathsetmacro{\x}{11.65} \pgfmathsetmacro{\y}{0.8}
      % \draw[lightgray, step=0.5] (0,0) grid (14,6);
      \draw[->, ultra thick, PastelGreen] (\x+1, \y) -- (\x, \y);
    \end{tikzpicture}%
  \end{frame}

  \subsection{Forward Pass}
  \begin{frame}{Forward Pass}
    \centering
    \vspace{5pt}
    \begin{tikzpicture}[remember picture]
    \pgfmathsetmacro{\ex}{1} \pgfmathsetmacro{\ey}{.9} \pgfmathsetmacro{\eradius}{0.9}

    % second
    \fill (-1, 0) circle[x radius=0.15, y radius=0.15];
    \node at (-1.5, 0) {$x_{1}$};

    \draw[thick, ->] (-1, 0) -- (\ex+8.5, 0);

    \node[font=\normalsize] (theta1) at (\ex-0.9, 0.3) {$\theta^{(1)}$};
    \node[font=\normalsize] (theta2) at (\ex+3.3, 0.3) {$\theta^{(2)}$};

    \draw[fill=PastelBlue] (\ex+1.5,\ey) arc (90:270:\eradius) -- cycle;
    \draw[fill=PastelBlue2] (\ex+1.5,\ey) arc (90:-90:\eradius) -- cycle;
    \node[font=\Large] (z1) at (\ex + 1.1, 0.1) {$z^{(1)}$};
    \node[font=\Large] (a1) at (\ex + 2, 0.1) {$a^{(1)}$};%

    \draw[fill=PastelBlue] (\ex+5.5,\ey) arc (90:270:\eradius) -- cycle;
    \draw[fill=PastelBlue2] (\ex+5.5,\ey) arc (90:-90:\eradius) -- cycle;
    \node[font=\Large] (z2) at (\ex + 5.1, 0.1) {$z^{(2)}$};
    \node[font=\Large] (a2) at (\ex + 6, 0.1) {$a^{(2)}$};%
    % \node[font=\scriptsize] at (\ex + 0.5, 0) {$0.1234$};
    % \node[font=\scriptsize] at (\ex + 1.5, 0) {$0.5672$};%
    % \node[font=\scriptsize] at (\ex + 4.5, 0) {$0.1234$};
    % \node[font=\scriptsize] at (\ex + 5.5, 0) {$0.1234$};
\end{tikzpicture}
    \vspace{10pt}

    \begin{itemize}
      \item Vamos considerar os pesos iniciais como: $\theta^{(1)}=-0.5$ e
        $\theta^{(2)}=0.5$.

      \item Computando o loss para o par $x_{1}=1$, $y=1$.
    \end{itemize}
    \vfill
    \textbf{Passo a passo (Forward Pass):}
    \begin{columns}
      \begin{column}{0.6\textwidth}
        \begin{enumerate}
          \visible<1->{\item
          $z^{(1)}= x_{1}\cdot \theta^{(1)}= 1 \cdot (-0.5) = -0.5$} \visible<2->{\item
          $a^{(1)}= \sigma(z^{(1)}) = \sigma(-0.5) = 0.3775$} \visible<3->{\item
          $z^{(2)}= a^{(1)}\cdot \theta^{(2)}= 0.3775 \cdot 0.5 = 0.18875$} \visible<4->{\item
          $\hat{y}= a^{(2)}= \sigma(z^{(2)}) = \sigma(0.18875) = 0.5470$}
        \end{enumerate}
      \end{column}
      \begin{column}{0.4\textwidth}
        \visible<4->{
        \begin{block}{Resumindo}
          Esse é o \textit{forward pass} da rede neural, onde calculamos as ativações
          e obtemos a saída da rede $\hat{y}= a^{(2)}= 0.5470$
        \end{block}}
      \end{column}
    \end{columns}
  \end{frame}
  \subsection{Função de Custo}

  \begin{frame}{Loss Function}
    \vspace{5pt}
    \begin{figure}
      \centering
      % Use the tikz file to draw the MLP
      \begin{tikzpicture}[remember picture]
    \pgfmathsetmacro{\ex}{1} \pgfmathsetmacro{\ey}{.9} \pgfmathsetmacro{\eradius}{0.9}

    % second
    \fill (-1, 0) circle[x radius=0.15, y radius=0.15];
    \node at (-1.5, 0) {$x_{1}$};

    \draw[thick, ->] (-1, 0) -- (\ex+8.5, 0);

    \node[font=\normalsize] (theta1) at (\ex-0.9, 0.3) {$\theta^{(1)}$};
    \node[font=\normalsize] (theta2) at (\ex+3.3, 0.3) {$\theta^{(2)}$};

    \draw[fill=PastelBlue] (\ex+1.5,\ey) arc (90:270:\eradius) -- cycle;
    \draw[fill=PastelBlue2] (\ex+1.5,\ey) arc (90:-90:\eradius) -- cycle;
    \node[font=\Large] (z1) at (\ex + 1.1, 0.1) {$z^{(1)}$};
    \node[font=\Large] (a1) at (\ex + 2, 0.1) {$a^{(1)}$};%

    \draw[fill=PastelBlue] (\ex+5.5,\ey) arc (90:270:\eradius) -- cycle;
    \draw[fill=PastelBlue2] (\ex+5.5,\ey) arc (90:-90:\eradius) -- cycle;
    \node[font=\Large] (z2) at (\ex + 5.1, 0.1) {$z^{(2)}$};
    \node[font=\Large] (a2) at (\ex + 6, 0.1) {$a^{(2)}$};%
    % \node[font=\scriptsize] at (\ex + 0.5, 0) {$0.1234$};
    % \node[font=\scriptsize] at (\ex + 1.5, 0) {$0.5672$};%
    % \node[font=\scriptsize] at (\ex + 4.5, 0) {$0.1234$};
    % \node[font=\scriptsize] at (\ex + 5.5, 0) {$0.1234$};
\end{tikzpicture}
    \end{figure}
    \vspace{10pt}
    \vfill
    Nessa etapa computamos a \textit{Loss Function} para avaliar o desempenho da
    rede.

    Loss: $J(\theta)=\frac{1}{2}(y-\hat{y})^{2}$.
    \[
      J(\theta) = \frac{1}{2}(1 - 0.5470)^{2}= 0.1026045
    \]
    \only<2>{
    \begin{block}{Resumindo}
      Essa é a etapa de computar o valor da \textit{Loss Function}. Observe que o
      valor da \textit{Loss Function} não será usado diretamente nos gradientes
      nem na atualização dos pesos
    \end{block}
    }
  \end{frame}

  \begin{frame}{Comparando com PyTorch}
    \begin{columns}
      \begin{column}{.5\linewidth}
        \vfill
        \begin{figure}
          \resizebox{\textwidth}{!}{
          \begin{tikzpicture}[remember picture]
    \pgfmathsetmacro{\ex}{1} \pgfmathsetmacro{\ey}{.9} \pgfmathsetmacro{\eradius}{0.9}

    % second
    \fill (-1, 0) circle[x radius=0.15, y radius=0.15];
    \node at (-1.5, 0) {$x_{1}$};

    \draw[thick, ->] (-1, 0) -- (\ex+8.5, 0);

    \node[font=\normalsize] (theta1) at (\ex-0.9, 0.3) {$\theta^{(1)}$};
    \node[font=\normalsize] (theta2) at (\ex+3.3, 0.3) {$\theta^{(2)}$};

    \draw[fill=PastelBlue] (\ex+1.5,\ey) arc (90:270:\eradius) -- cycle;
    \draw[fill=PastelBlue2] (\ex+1.5,\ey) arc (90:-90:\eradius) -- cycle;
    \node[font=\Large] (z1) at (\ex + 1.1, 0.1) {$z^{(1)}$};
    \node[font=\Large] (a1) at (\ex + 2, 0.1) {$a^{(1)}$};%

    \draw[fill=PastelBlue] (\ex+5.5,\ey) arc (90:270:\eradius) -- cycle;
    \draw[fill=PastelBlue2] (\ex+5.5,\ey) arc (90:-90:\eradius) -- cycle;
    \node[font=\Large] (z2) at (\ex + 5.1, 0.1) {$z^{(2)}$};
    \node[font=\Large] (a2) at (\ex + 6, 0.1) {$a^{(2)}$};%
    % \node[font=\scriptsize] at (\ex + 0.5, 0) {$0.1234$};
    % \node[font=\scriptsize] at (\ex + 1.5, 0) {$0.5672$};%
    % \node[font=\scriptsize] at (\ex + 4.5, 0) {$0.1234$};
    % \node[font=\scriptsize] at (\ex + 5.5, 0) {$0.1234$};
\end{tikzpicture}
          }
        \end{figure}
        \vfill
        \begin{align*}
          \hat{y}   & = 0.5470    \\
          J(\theta) & = 0.1026045
        \end{align*}
        \vfill
      \end{column}
      \begin{column}{.5\linewidth}
        \centering
        \resizebox*{!}{.75\textheight}{ \lstinputlisting[style=vscode-light]{slides/code/01_forward_pass.py}%
        }
      \end{column}
    \end{columns}
  \end{frame}

  \subsection{Backward Pass}

  \begin{frame}{Backward Pass}
    \vspace{5pt}
    \begin{figure}
      \centering
      % Use the tikz file to draw the MLP
      \begin{tikzpicture}[remember picture]
    \pgfmathsetmacro{\ex}{1} \pgfmathsetmacro{\ey}{.9} \pgfmathsetmacro{\eradius}{0.9}

    % second
    \fill (-1, 0) circle[x radius=0.15, y radius=0.15];
    \node at (-1.5, 0) {$x_{1}$};

    \draw[thick, ->] (-1, 0) -- (\ex+8.5, 0);

    \node[font=\normalsize] (theta1) at (\ex-0.9, 0.3) {$\theta^{(1)}$};
    \node[font=\normalsize] (theta2) at (\ex+3.3, 0.3) {$\theta^{(2)}$};

    \draw[fill=PastelBlue] (\ex+1.5,\ey) arc (90:270:\eradius) -- cycle;
    \draw[fill=PastelBlue2] (\ex+1.5,\ey) arc (90:-90:\eradius) -- cycle;
    \node[font=\Large] (z1) at (\ex + 1.1, 0.1) {$z^{(1)}$};
    \node[font=\Large] (a1) at (\ex + 2, 0.1) {$a^{(1)}$};%

    \draw[fill=PastelBlue] (\ex+5.5,\ey) arc (90:270:\eradius) -- cycle;
    \draw[fill=PastelBlue2] (\ex+5.5,\ey) arc (90:-90:\eradius) -- cycle;
    \node[font=\Large] (z2) at (\ex + 5.1, 0.1) {$z^{(2)}$};
    \node[font=\Large] (a2) at (\ex + 6, 0.1) {$a^{(2)}$};%
    % \node[font=\scriptsize] at (\ex + 0.5, 0) {$0.1234$};
    % \node[font=\scriptsize] at (\ex + 1.5, 0) {$0.5672$};%
    % \node[font=\scriptsize] at (\ex + 4.5, 0) {$0.1234$};
    % \node[font=\scriptsize] at (\ex + 5.5, 0) {$0.1234$};
\end{tikzpicture}
    \end{figure}
    \vspace{10pt}
    \begin{columns}
      \begin{column}{0.6\textwidth}
        Nossa rede neural computa a seguinte função:
        \[
          a^{(2)}= \sigma(\sigma(x_{1}\cdot \theta^{(1)})\cdot \theta^{(2)})
        \]

        Logo, a função de custo é:
        \[
          J(\theta)=\frac{1}{2}(y-a^{(2)})^{2}=\frac{1}{2}\left(1-\sigma\left(\sigma
          ( x_{1}\cdot \theta^{(1)})\cdot \theta^{(2)}\right)\right)^{2}
        \]
      \end{column}
      \begin{column}{0.4\textwidth}
        \vfill
        \visible<2>{\begin{alertblock}{Problema}
          Como computar o gradiente da função de custo em relação aos pesos $\theta
          ^{(1)}$ e $\theta^{(2)}$?
          \begin{align*}
            \frac{\partial J}{\partial\theta^{(2)}} & = ? \\
            \frac{\partial J}{\partial\theta^{(1)}} & = ?
          \end{align*}
        \end{alertblock}}
      \end{column}
    \end{columns}
  \end{frame}

  \begin{frame}{Backward Pass}
    \vspace{5pt}
    \begin{figure}
      \centering
      % Use the tikz file to draw the MLP
      \begin{tikzpicture}[remember picture]
    \pgfmathsetmacro{\ex}{1} \pgfmathsetmacro{\ey}{.9} \pgfmathsetmacro{\eradius}{0.9}

    % second
    \fill (-1, 0) circle[x radius=0.15, y radius=0.15];
    \node at (-1.5, 0) {$x_{1}$};

    \draw[thick, ->] (-1, 0) -- (\ex+8.5, 0);

    \node[font=\normalsize] (theta1) at (\ex-0.9, 0.3) {$\theta^{(1)}$};
    \node[font=\normalsize] (theta2) at (\ex+3.3, 0.3) {$\theta^{(2)}$};

    \draw[fill=PastelBlue] (\ex+1.5,\ey) arc (90:270:\eradius) -- cycle;
    \draw[fill=PastelBlue2] (\ex+1.5,\ey) arc (90:-90:\eradius) -- cycle;
    \node[font=\Large] (z1) at (\ex + 1.1, 0.1) {$z^{(1)}$};
    \node[font=\Large] (a1) at (\ex + 2, 0.1) {$a^{(1)}$};%

    \draw[fill=PastelBlue] (\ex+5.5,\ey) arc (90:270:\eradius) -- cycle;
    \draw[fill=PastelBlue2] (\ex+5.5,\ey) arc (90:-90:\eradius) -- cycle;
    \node[font=\Large] (z2) at (\ex + 5.1, 0.1) {$z^{(2)}$};
    \node[font=\Large] (a2) at (\ex + 6, 0.1) {$a^{(2)}$};%
    % \node[font=\scriptsize] at (\ex + 0.5, 0) {$0.1234$};
    % \node[font=\scriptsize] at (\ex + 1.5, 0) {$0.5672$};%
    % \node[font=\scriptsize] at (\ex + 4.5, 0) {$0.1234$};
    % \node[font=\scriptsize] at (\ex + 5.5, 0) {$0.1234$};
\end{tikzpicture}
    \end{figure}
    \vspace{10pt}

    \textbf{Solução:} A função $J(\theta)$ é uma função composta, então usamos a
    regra da cadeia
    \vfill
    \textbf{Gradiente para $\theta^{(2)}$:}
    \[
      \frac{\partial J}{\partial\theta^{(2)}}=\frac{\partial J}{\partial a^{(2)}}
      \frac{\partial a^{(2)}}{\partial z^{(2)}}\frac{\partial z^{(2)}}{\partial\theta^{(2)}}
    \]
    \vfill
    \textbf{Gradiente para $\theta^{(1)}$:}
    \[
      \frac{\partial J}{\partial\theta^{(1)}}=\frac{\partial J}{\partial a^{(2)}}
      \frac{\partial a^{(2)}}{\partial z^{(2)}}\frac{\partial z^{(2)}}{\partial
      a^{(1)}}\frac{\partial a^{(1)}}{\partial z^{(1)}}\frac{\partial z^{(1)}}{\partial\theta^{(1)}}
    \]
  \end{frame}

  \begin{frame}{Backward Pass}
    % \vspace{5pt}
    %     \begin{figure}
    %       \centering
    %       % Use the tikz file to draw the MLP
    %       \begin{tikzpicture}[remember picture]
    \pgfmathsetmacro{\ex}{1} \pgfmathsetmacro{\ey}{.9} \pgfmathsetmacro{\eradius}{0.9}

    % second
    \fill (-1, 0) circle[x radius=0.15, y radius=0.15];
    \node at (-1.5, 0) {$x_{1}$};

    \draw[thick, ->] (-1, 0) -- (\ex+8.5, 0);

    \node[font=\normalsize] (theta1) at (\ex-0.9, 0.3) {$\theta^{(1)}$};
    \node[font=\normalsize] (theta2) at (\ex+3.3, 0.3) {$\theta^{(2)}$};

    \draw[fill=PastelBlue] (\ex+1.5,\ey) arc (90:270:\eradius) -- cycle;
    \draw[fill=PastelBlue2] (\ex+1.5,\ey) arc (90:-90:\eradius) -- cycle;
    \node[font=\Large] (z1) at (\ex + 1.1, 0.1) {$z^{(1)}$};
    \node[font=\Large] (a1) at (\ex + 2, 0.1) {$a^{(1)}$};%

    \draw[fill=PastelBlue] (\ex+5.5,\ey) arc (90:270:\eradius) -- cycle;
    \draw[fill=PastelBlue2] (\ex+5.5,\ey) arc (90:-90:\eradius) -- cycle;
    \node[font=\Large] (z2) at (\ex + 5.1, 0.1) {$z^{(2)}$};
    \node[font=\Large] (a2) at (\ex + 6, 0.1) {$a^{(2)}$};%
    % \node[font=\scriptsize] at (\ex + 0.5, 0) {$0.1234$};
    % \node[font=\scriptsize] at (\ex + 1.5, 0) {$0.5672$};%
    % \node[font=\scriptsize] at (\ex + 4.5, 0) {$0.1234$};
    % \node[font=\scriptsize] at (\ex + 5.5, 0) {$0.1234$};
\end{tikzpicture}
    %     \end{figure}
    %     \vspace{10pt}
    \[
      \frac{\partial J}{\partial\theta^{(2)}}=\frac{\partial J}{\partial a^{(2)}}
      \frac{\partial a^{(2)}}{\partial z^{(2)}}\frac{\partial z^{(2)}}{\partial\theta^{(2)}}
    \]
    \begin{columns}
      \begin{column}{.33\linewidth}
        \begin{itemize}
          \item<1-> $\frac{\partial J}{\partial a^{(2)}}=a^{(2)}-y$

          \item<2-> $\frac{\partial a^{(2)}}{\partial z^{(2)}}=a^{(2)}(1-a^{(2)})$

          \item<3-> $\frac{\partial z^{(2)}}{\partial\theta^{(2)}}=a^{(1)}$
        \end{itemize}
      \end{column}
      \begin{column}{.67\linewidth}
        \only<1>{
        \begin{align*}
          \frac{\partial J}{\partial a^{(2)}} & =\frac{\partial \left( \frac{1}{2}(y-a^{(2)})^{2}\right)}{\partial a^{(2)}} \\
                                              & =(-1)\cancel{2}\frac{1}{\cancel{2}}(y-a^{(2)})                              \\
                                              & =a^{(2)}-y
        \end{align*}} \only<2>{
        \begin{align*}
          \frac{\partial a^{(2)}}{\partial z^{(2)}} & =\frac{\partial \left( \frac{1}{1+e^{-z^{(2)}}}\right)}{\partial z^{(2)}}=\frac{\partial \left( (1+e^{-z^{(2)}})^{-1}\right)}{\partial z^{(2)}}                       \\
                                                    & =(\cancel{-1})(1+e^{-z^{(2)}})^{-2}e^{-z^{(2)}}(\cancel{-1})                                                                                                          \\
                                                    & =\frac{\textcolor{red}{1}+e^{-z^{(2)}}\textcolor{red}{-1}}{(1+e^{-z^{(2)}})^{2}}                                                                                      \\
                                                    & =\frac{1}{1+e^{-z^{(2)}}}\left(\cancelto{1}{\frac{\textcolor{red}{1}+e^{-z^{(2)}}}{1+e^{-z^{(2)}}}}\textcolor{red}{-}\frac{\textcolor{red}{1}}{1+e^{-z^{(2)}}}\right) \\
                                                    & =\sigma(z^{(2)})(1-\sigma(z^{(2)})) = a^{(2)}(1-a^{(2)})
        \end{align*}} \only<3>{
        \begin{align*}
          \frac{\partial z^{(2)}}{\partial \theta^{(2)}} & =\frac{\partial \left(a^{(1)}\theta^{(2)}\right)}{\partial \theta^{(2)}} \\
                                                         & =a^{(1)}
        \end{align*}}
      \end{column}
    \end{columns}
  \end{frame}

  \begin{frame}{Backward Pass}
    % \vspace{5pt}
    %     \begin{figure}
    %       \centering
    %       % Use the tikz file to draw the MLP
    %       \begin{tikzpicture}[remember picture]
    \pgfmathsetmacro{\ex}{1} \pgfmathsetmacro{\ey}{.9} \pgfmathsetmacro{\eradius}{0.9}

    % second
    \fill (-1, 0) circle[x radius=0.15, y radius=0.15];
    \node at (-1.5, 0) {$x_{1}$};

    \draw[thick, ->] (-1, 0) -- (\ex+8.5, 0);

    \node[font=\normalsize] (theta1) at (\ex-0.9, 0.3) {$\theta^{(1)}$};
    \node[font=\normalsize] (theta2) at (\ex+3.3, 0.3) {$\theta^{(2)}$};

    \draw[fill=PastelBlue] (\ex+1.5,\ey) arc (90:270:\eradius) -- cycle;
    \draw[fill=PastelBlue2] (\ex+1.5,\ey) arc (90:-90:\eradius) -- cycle;
    \node[font=\Large] (z1) at (\ex + 1.1, 0.1) {$z^{(1)}$};
    \node[font=\Large] (a1) at (\ex + 2, 0.1) {$a^{(1)}$};%

    \draw[fill=PastelBlue] (\ex+5.5,\ey) arc (90:270:\eradius) -- cycle;
    \draw[fill=PastelBlue2] (\ex+5.5,\ey) arc (90:-90:\eradius) -- cycle;
    \node[font=\Large] (z2) at (\ex + 5.1, 0.1) {$z^{(2)}$};
    \node[font=\Large] (a2) at (\ex + 6, 0.1) {$a^{(2)}$};%
    % \node[font=\scriptsize] at (\ex + 0.5, 0) {$0.1234$};
    % \node[font=\scriptsize] at (\ex + 1.5, 0) {$0.5672$};%
    % \node[font=\scriptsize] at (\ex + 4.5, 0) {$0.1234$};
    % \node[font=\scriptsize] at (\ex + 5.5, 0) {$0.1234$};
\end{tikzpicture}
    %     \end{figure}
    %     \vspace{10pt}
    \[
      \frac{\partial J}{\partial\theta^{(1)}}= \frac{\partial J}{\partial a^{(2)}}
      \frac{\partial a^{(2)}}{\partial z^{(2)}}\frac{\partial z^{(2)}}{\partial
      a^{(1)}}\frac{\partial a^{(1)}}{\partial z^{(1)}}\frac{\partial z^{(1)}}{\partial\theta^{(1)}}
    \]
    \begin{columns}
      \begin{column}{.33\linewidth}
        \begin{itemize}
          \item $\frac{\partial J}{\partial a^{(2)}}=a^{(2)}-y$

          \item $\frac{\partial a^{(2)}}{\partial z^{(2)}}=a^{(2)}(1-a^{(2)})$

          \item<2-> $\frac{\partial z^{(2)}}{\partial a^{(1)}}=\theta^{(2)}$

          \item<3-> $\frac{\partial a^{(1)}}{\partial z^{(1)}}=a^{(1)}(1-a^{(1)})$

          \item<4> $\frac{\partial z^{(1)}}{\partial \theta^{(1)}}=x_{1}$
        \end{itemize}
      \end{column}
      \begin{column}{.67\linewidth}
        \only<1>{ Os primeiros termos são iguais aos dois primeiros termos do gradiente
        de $\theta^{(2)}$

        } \only<2>{
        \begin{align*}
          \frac{\partial z^{(2)}}{\partial a^{(1)}} & =\frac{\partial \left(a^{(1)}\theta^{(2)}\right)}{\partial a^{(1)}} \\
                                                    & =\theta^{(2)}
        \end{align*}
        } \only<3>{
        \begin{align*}
          \frac{\partial a^{(1)}}{\partial z^{(1)}} & =\frac{\partial \left( \frac{1}{1+e^{-z^{(1)}}}\right)}{\partial z^{(1)}}=\frac{\partial \left( (1+e^{-z^{(1)}})^{-1}\right)}{\partial z^{(1)}}                       \\
                                                    & =(\cancel{-1})(1+e^{-z^{(1)}})^{-2}e^{-z^{(1)}}(\cancel{-1})                                                                                                          \\
                                                    & =\frac{\textcolor{red}{1}+e^{-z^{(1)}}\textcolor{red}{-1}}{(1+e^{-z^{(1)}})^{2}}                                                                                      \\
                                                    & =\frac{1}{1+e^{-z^{(1)}}}\left(\cancelto{1}{\frac{\textcolor{red}{1}+e^{-z^{(1)}}}{1+e^{-z^{(1)}}}}\textcolor{red}{-}\frac{\textcolor{red}{1}}{1+e^{-z^{(1)}}}\right) \\
                                                    & =\sigma(z^{(1)})(1-\sigma(z^{(1)}))=a^{(1)}(1-a^{(1)})
        \end{align*}} \only<4>{
        \begin{align*}
          \frac{\partial z^{(1)}}{\partial \theta^{(1)}} & =\frac{\partial \left(x_{1}\theta^{(1)}\right)}{\partial \theta^{(1)}} \\
                                                         & =x_{1}
        \end{align*}}
      \end{column}
    \end{columns}
  \end{frame}

  \begin{frame}{Backward Pass}
    \textbf{Vamos Computar o Gradiente}
    \begin{align*}
      \onslide<1->{\frac{\partial J}{\partial\theta^{(2)}} & = \underbrace{(a^{(2)}-y) \cdot a^{(2)}(1 -a^{(2)})}_{\delta^{(2)}}\cdot a^{(1)}}                                          \\
      \onslide<2->{                                            & = (0.5470-1) \cdot 0.5470(1-0.5470) \cdot 0.3775 = \textbf{−0.04237412}}                                                   \\
      \onslide<3->{\frac{\partial J}{\partial\theta^{(1)}} & = \underbrace{(a^{(2)}-y) \cdot a^{(2)}(1 -a^{(2)})}_{\delta^{(2)}}\cdot \theta^{(2)}\cdot a^{(1)}(1-a^{(1)}) \cdot x_{1}} \\
      \onslide<4->{                                            & = (0.5470-1) \cdot 0.5470(1-0.5470) \cdot 0.5 \cdot 0.3775(1-0.3775) \cdot 1 = \textbf{−0.01318894}}
    \end{align*}
    \only<5>{
    \begin{block}{Resumindo}
      No \textit{Backward Pass} calculamos os gradientes da função de custo em
      relação a todos os parâmetros treináveis
    \end{block}
    }
  \end{frame}

  \begin{frame}{Descida de Gradiente}
    \begin{itemize}
      \item Usando $\eta = 0.1$: $\theta_{t}= \theta_{t-1}- \eta \nabla_{\theta} J$

      \item $\theta_{t}^{(2)}= 0.5 - 0.1(-0.042374) = \textbf{0.5042374}$

      \item $\theta_{t}^{(1)}= -0.5 - 0.1(-0.013189) = \textbf{-0.4986811}$
    \end{itemize}
    \vfill
    \only<2>{
    \begin{block}{Resumindo}
      Na \textit{Descida de Gradiente} usamos o gradiente para atualizar os
      parâmetros treináveis
    \end{block}
    }
  \end{frame}

  \begin{frame}{Resultado da Atualização}
    \begin{itemize}
      \item Após a atualização, os pesos da rede são:
        \begin{itemize}
          \item $\theta^{(1)}=-0.4986811$

          \item $\theta^{(2)}=0.5042374$
        \end{itemize}

      \item O novo valor de $a^{(2)}$ é:
        \[
          a^{(2)}= \sigma(\sigma(x_{1}\cdot\theta^{(1)})\cdot\theta^{(2)})= \sigma
          (\sigma(1\cdot(-0.4986811))\cdot 0.5042374)=0.547\textcolor{PastelBlue2}{4}
        \]

      \item O novo valor de $J(\theta)$ é:
        \[
          J(\theta)=\frac{1}{2}(y-a^{(2)})^{2}=\frac{1}{2}(1-0.5474)^{2}=0.102\textcolor
          {PastelBlue2}{4234}
        \]
        \visible<2>{\item O novo valor de $J(\theta)$ é menor que o anterior
        $0.102\textcolor{Red}{6045}$, indicando que o erro do MLP é menor}
    \end{itemize}
    \vfill
  \end{frame}

  \subsection{Descida de Gradiente}

  \begin{frame}[fragile,containsverbatim]
    \frametitle{Comparando com PyTorch}
    \begin{columns}
      \begin{column}{0.4\textwidth}
        \begin{align*}
          \frac{\partial J}{\partial \mathbf{\Theta}^{(1)}} & =-0.01318894 \\
          \frac{\partial J}{\partial \mathbf{\Theta}^{(2)}} & =-0.04237412 \\
          \mathbf{\Theta}_{t}^{(2)}                         & = 0.5042374  \\
          \mathbf{\Theta}_{t}^{(1)}                         & = -0.4986811
        \end{align*}
      \end{column}
      \begin{column}{.6\linewidth}
        \centering
        \resizebox*{\linewidth}{!}{ \lstinputlisting[style=vscode-light]{slides/code/01_backward_pass.py}%
        }
      \end{column}
    \end{columns}
  \end{frame}

  \section{Vetorização do Treinamento}

  \begin{frame}{Generalizando o Treinamento}
  \begin{figure}
      \centering
      \begin{tikzpicture}[remember picture]
    \pgfmathsetmacro{\ex}{1} \pgfmathsetmacro{\ey}{.9} \pgfmathsetmacro{\eradius}{0.9}

    % second
    \fill (-1, 0) circle[x radius=0.15, y radius=0.15];
    \node at (-1.5, 0) {$x_{1}$};

    \draw[thick, ->] (-1, 0) -- (\ex+8.5, 0);

    \node[font=\normalsize] (theta1) at (\ex-0.9, 0.3) {$\theta^{(1)}$};
    \node[font=\normalsize] (theta2) at (\ex+3.3, 0.3) {$\theta^{(2)}$};

    \draw[fill=PastelBlue] (\ex+1.5,\ey) arc (90:270:\eradius) -- cycle;
    \draw[fill=PastelBlue2] (\ex+1.5,\ey) arc (90:-90:\eradius) -- cycle;
    \node[font=\Large] (z1) at (\ex + 1.1, 0.1) {$z^{(1)}$};
    \node[font=\Large] (a1) at (\ex + 2, 0.1) {$a^{(1)}$};%

    \draw[fill=PastelBlue] (\ex+5.5,\ey) arc (90:270:\eradius) -- cycle;
    \draw[fill=PastelBlue2] (\ex+5.5,\ey) arc (90:-90:\eradius) -- cycle;
    \node[font=\Large] (z2) at (\ex + 5.1, 0.1) {$z^{(2)}$};
    \node[font=\Large] (a2) at (\ex + 6, 0.1) {$a^{(2)}$};%
    % \node[font=\scriptsize] at (\ex + 0.5, 0) {$0.1234$};
    % \node[font=\scriptsize] at (\ex + 1.5, 0) {$0.5672$};%
    % \node[font=\scriptsize] at (\ex + 4.5, 0) {$0.1234$};
    % \node[font=\scriptsize] at (\ex + 5.5, 0) {$0.1234$};
\end{tikzpicture}
  \end{figure}
  \begin{columns}
      \begin{column}{.55\linewidth}
        Repare que:
          \begin{itemize}
              \visible<1->{\item  Ambas as derivadas são computadas com um termo $\delta^{(l)}$ multiplicado pela entrada da camada $\frac{\partial z^{(l)}}{\partial \theta^{(l)}}$}
              \visible<2>{\item $\delta^{(1)}$ pode ser obtido a partir de $\delta^{(2)}$}
          \end{itemize}
      \end{column}
      \begin{column}{.45\linewidth}
        \begin{align*}
          \frac{\partial J}{\partial\theta^{(2)}}&=\underbrace{\textcolor{red}{\frac{\partial J}{\partial a^{(2)}}
          \frac{\partial a^{(2)}}{\partial z^{(2)}}}}_{\textcolor{red}{\delta^{(2)}}}\textcolor{teal}{\frac{\partial z^{(2)}}{\partial\theta^{(2)}}}\\
          \frac{\partial J}{\partial\theta^{(1)}}&=\underbrace{\underbrace{\textcolor{red}{\frac{\partial J}{\partial a^{(2)}}
          \frac{\partial a^{(2)}}{\partial z^{(2)}}}}_{\textcolor{red}{\delta^{(2)}}}\textcolor{magenta}{\frac{\partial z^{(2)}}{\partial
          a^{(1)}}\frac{\partial a^{(1)}}{\partial z^{(1)}}}}_{\textcolor{magenta}{\delta^{(1)}}}\textcolor{blue}{\frac{\partial z^{(1)}}{\partial\theta^{(1)}}}
      \end{align*}        
      \end{column}
  \end{columns}
  
  \begin{tikzpicture}[overlay, remember picture]
      % \draw[lightgray, step=0.5] (0,0) grid (14,7);
      \draw[ultra thick, ->, teal] ($(theta2)-(1.cm,0)$) to (theta2);
      \draw[ultra thick, ->, red] ($(theta2)+(1.cm,0)$) to (theta2);
      
      \draw[ultra thick, ->, blue] ($(theta1)-(1.cm,0)$) to (theta1);
      \draw[ultra thick, ->, magenta] ($(theta1)+(1.cm,0)$) to (theta1);
    \end{tikzpicture}
  \end{frame}


  
  \begin{frame}{Generalizando o Treinamento}
    O gradiente em cada camada pode ser dividido em dois componentes:
    \begin{itemize}
      \item \textbf{Delta da Camada:}
        $\bm{\delta}^{(l)}= \frac{\partial J}{\partial \mathbf{z}^{(l)}}$

      \item \textbf{Ativação da Camada Anterior:} $\mathbf{a}^{(l-1)}$
    \end{itemize}
    \vfill
    \textbf{Backpropagation:}
    \begin{enumerate}
      \item \textbf{Erro na Camada de Saída:}
        $\bm{\delta}^{(L)}=\frac{\partial J}{\partial \mathbf{a}^{(L)}}\odot\frac{\partial \mathbf{a}^{(L)}}{\partial \mathbf{z}^{(L)}}$
        \vfill

      \item \textbf{Propagação do Erro:}
        $\bm{\delta}^{(l)}=(\bm{\delta}^{(l+1)}\mathbf{\Theta}^{(l+1)})\odot \mathbf{a}^{\prime(l)}$
        \vfill

      \item \textbf{Gradiente do Parâmetros:}
        $\frac{\partial J}{\partial\mathbf{\Theta}^{(i)}}=\bm{\delta}^{(i)^\top}\mathbf{a}^{(i-1)}$
        \vfill


    \end{enumerate}
  \end{frame}

  \begin{frame}{Exemplo Vetorizado}
    \begin{columns}
        \begin{column}{.6\textwidth}

        \begin{itemize}
            \item Dataset (\textit{Batch Size} = 2)
            \[
              \mathbf{X} =
              \begin{bmatrix}
                0.5 & 0.1 \\
                0.2 & 0.6
              \end{bmatrix}, \quad \mathbf{y} =
              \begin{bmatrix}
                0.6 \\
                0.8
              \end{bmatrix}
            \]
            \item Arquitetura
                \begin{itemize}
                  \item Camada de Entrada: 2 unidades
                  \item Camada Oculta 1: 3 unidades
                  \item Camada de Saída: 1 unidade
                \end{itemize}
            \item Pesos Iniciais
                \begin{itemize}
                  \item Camada 1: $\mathbf{\Theta}^{(1)}=
                    \begin{bmatrix}
                      0.5  & 0.2  \\
                      0.6  & -0.1 \\
                      -0.4 & -0.3
                    \end{bmatrix}$
            
                  \item Camada 2: $\mathbf{\Theta}^{(2)}=
                    \begin{bmatrix}
                      0.7 & -0.1 & 0.2
                    \end{bmatrix}$
                \end{itemize}
        \end{itemize}

        \end{column}
        \begin{column}{.4\textwidth}
            \begin{figure}
                \centering
                \resizebox{\linewidth}{!}{
                \begin{tikzpicture}[remember picture]
    \pgfmathsetmacro{\ex}{1} \pgfmathsetmacro{\ey}{.9} \pgfmathsetmacro{\eradius}{0.9}
    
    \fill (-1, 1) circle[x radius=0.15, y radius=0.15];
    \node at (-1.5, 1) {$x_{1}$};

    \draw[thick, ->] (-1, 1) -- (\ex+1., 0) node[pos=0.6, above]  {$\theta^{(1)}_{21}$};
    \draw[thick, ->] (-1, 1) -- (\ex+1., 2.5) node[pos=0.6, above]  {$\theta^{(1)}_{11}$};
    \draw[thick, ->] (-1, 1) -- (\ex+1., -2.5) node[pos=0.6, above]  {$\theta^{(1)}_{31}$};
    
    \fill (-1, -1) circle[x radius=0.15, y radius=0.15];
    \node at (-1.5, -1) {$x_{2}$};
    \draw[thick, ->] (-1, -1) -- (\ex+1., 2.5) node[pos=0.6, above] {$\theta^{(1)}_{12}$};
    \draw[thick, ->] (-1, -1) -- (\ex+1., 0) node[pos=0.6,above] {$\theta^{(1)}_{22}$};
    \draw[thick, ->] (-1, -1) -- (\ex+1., -2.5) node[pos=0.6, above] {$\theta^{(1)}_{32}$};

    \draw[thick, ->] (\ex+1., 2) -- (\ex+4.5, 0) node[pos=0.41, above]  {$\theta^{(2)}_{11}$};
    \draw[thick, ->] (\ex+1., 0) -- (\ex+6, 0) node[pos=0.3, above]  {$\theta^{(2)}_{12}$};
    \draw[thick, ->] (\ex+1., -2) -- (\ex+4.5, 0) node[pos=0.41, above]  {$\theta^{(2)}_{13}$};

    \draw[fill=PastelBlue] (\ex+1,\ey+2) arc (90:270:\eradius) -- cycle;
    \draw[fill=PastelBlue2] (\ex+1,\ey+2) arc (90:-90:\eradius) -- cycle;
    \node[font=\Large] (z1_1) at (\ex + .6, 0.1+2) {$z^{(1)}_1$};
    \node[font=\Large] (a1_1) at (\ex + 1.5, 0.1+2.) {$a^{(1)}_1$};

    \draw[fill=PastelBlue] (\ex+1,\ey) arc (90:270:\eradius) -- cycle;
    \draw[fill=PastelBlue2] (\ex+1,\ey) arc (90:-90:\eradius) -- cycle;
    \node[font=\Large] (z1_2) at (\ex + .6, 0.1) {$z^{(1)}_2$};
    \node[font=\Large] (a1_2) at (\ex + 1.5, 0.1) {$a^{(1)}_2$};

    \draw[fill=PastelBlue] (\ex+1,\ey-2) arc (90:270:\eradius) -- cycle;
    \draw[fill=PastelBlue2] (\ex+1,\ey-2) arc (90:-90:\eradius) -- cycle;
    \node[font=\Large] (z1_3) at (\ex + .6, 0.1-2) {$z^{(1)}_3$};
    \node[font=\Large] (a1_3) at (\ex + 1.5, 0.1-2.) {$a^{(1)}_3$};

    \draw[fill=PastelBlue] (\ex+4,\ey) arc (90:270:\eradius) -- cycle;
    \draw[fill=PastelBlue2] (\ex+4,\ey) arc (90:-90:\eradius) -- cycle;
    \node[font=\Large] (z2) at (\ex + 3.6, 0.1) {$z^{(2)}_1$};
    \node[font=\Large] (a2) at (\ex + 4.5, 0.1) {$a^{(2)}_1$};
\end{tikzpicture}}
            \end{figure}
        \end{column}
    \end{columns}

  \end{frame}


  \begin{frame}
    \frametitle{Exemplo Vetorizado: Forward Pass}
    \begin{columns}
        \begin{column}{.6\linewidth}
        \small
        \begin{align*}
            \mathbf{Z}^{(1)} &= \mathbf{X} \cdot{\mathbf{\Theta}^{(1)}}^{\top}\\
            &= \begin{bmatrix}
        0.5 & 0.1 \\
        0.2 & 0.6
      \end{bmatrix}_{2\times2}\cdot
      \begin{bmatrix}
        0.5 & 0.6  & -0.4 \\
        0.2 & -0.1 & -0.3
      \end{bmatrix}_{2\times3}\\
      &=\begin{bmatrix}
        0.27 & 0.29 & -0.23 \\
        0.22 & 0.06 & -0.26
      \end{bmatrix}_{2\times3}\\
      \mathbf{A}^{(1)} &= \sigma(\mathbf{Z}^{(1)})\\
      &= \sigma\left(\begin{bmatrix}
        0.27 & 0.29 & -0.23 \\
        0.22 & 0.06 & -0.26
      \end{bmatrix}_{2\times3}\right)
      \\
      &=\begin{bmatrix}
        0.5671 & 0.5720 & 0.4428 \\
        0.5548 & 0.5150 & 0.4354
      \end{bmatrix}_{2\times3}
    \end{align*}                
    
        \end{column}
        \begin{column}{.4\textwidth}
            \begin{figure}
                \centering
                \resizebox{\linewidth}{!}{
                \begin{tikzpicture}[remember picture]
    \pgfmathsetmacro{\ex}{1} \pgfmathsetmacro{\ey}{.9} \pgfmathsetmacro{\eradius}{0.9}
    
    \fill (-1, 1) circle[x radius=0.15, y radius=0.15];
    \node at (-1.5, 1) {$x_{1}$};

    \draw[thick, ->] (-1, 1) -- (\ex+1., 0) node[pos=0.6, above]  {$\theta^{(1)}_{21}$};
    \draw[thick, ->] (-1, 1) -- (\ex+1., 2.5) node[pos=0.6, above]  {$\theta^{(1)}_{11}$};
    \draw[thick, ->] (-1, 1) -- (\ex+1., -2.5) node[pos=0.6, above]  {$\theta^{(1)}_{31}$};
    
    \fill (-1, -1) circle[x radius=0.15, y radius=0.15];
    \node at (-1.5, -1) {$x_{2}$};
    \draw[thick, ->] (-1, -1) -- (\ex+1., 2.5) node[pos=0.6, above] {$\theta^{(1)}_{12}$};
    \draw[thick, ->] (-1, -1) -- (\ex+1., 0) node[pos=0.6,above] {$\theta^{(1)}_{22}$};
    \draw[thick, ->] (-1, -1) -- (\ex+1., -2.5) node[pos=0.6, above] {$\theta^{(1)}_{32}$};

    \draw[thick, ->] (\ex+1., 2) -- (\ex+4.5, 0) node[pos=0.41, above]  {$\theta^{(2)}_{11}$};
    \draw[thick, ->] (\ex+1., 0) -- (\ex+6, 0) node[pos=0.3, above]  {$\theta^{(2)}_{12}$};
    \draw[thick, ->] (\ex+1., -2) -- (\ex+4.5, 0) node[pos=0.41, above]  {$\theta^{(2)}_{13}$};

    \draw[fill=PastelBlue] (\ex+1,\ey+2) arc (90:270:\eradius) -- cycle;
    \draw[fill=PastelBlue2] (\ex+1,\ey+2) arc (90:-90:\eradius) -- cycle;
    \node[font=\Large] (z1_1) at (\ex + .6, 0.1+2) {$z^{(1)}_1$};
    \node[font=\Large] (a1_1) at (\ex + 1.5, 0.1+2.) {$a^{(1)}_1$};

    \draw[fill=PastelBlue] (\ex+1,\ey) arc (90:270:\eradius) -- cycle;
    \draw[fill=PastelBlue2] (\ex+1,\ey) arc (90:-90:\eradius) -- cycle;
    \node[font=\Large] (z1_2) at (\ex + .6, 0.1) {$z^{(1)}_2$};
    \node[font=\Large] (a1_2) at (\ex + 1.5, 0.1) {$a^{(1)}_2$};

    \draw[fill=PastelBlue] (\ex+1,\ey-2) arc (90:270:\eradius) -- cycle;
    \draw[fill=PastelBlue2] (\ex+1,\ey-2) arc (90:-90:\eradius) -- cycle;
    \node[font=\Large] (z1_3) at (\ex + .6, 0.1-2) {$z^{(1)}_3$};
    \node[font=\Large] (a1_3) at (\ex + 1.5, 0.1-2.) {$a^{(1)}_3$};

    \draw[fill=PastelBlue] (\ex+4,\ey) arc (90:270:\eradius) -- cycle;
    \draw[fill=PastelBlue2] (\ex+4,\ey) arc (90:-90:\eradius) -- cycle;
    \node[font=\Large] (z2) at (\ex + 3.6, 0.1) {$z^{(2)}_1$};
    \node[font=\Large] (a2) at (\ex + 4.5, 0.1) {$a^{(2)}_1$};
\end{tikzpicture}}
            \end{figure}
        \end{column}    
    \end{columns}
  \end{frame}

\begin{frame}{Exemplo Vetorizado: Forward Pass}
    \begin{columns}
        \begin{column}{.6\linewidth}
        \small
        \begin{align*}
      \mathbf{Z}^{(2)}&= \mathbf{A}^{(1)}\cdot{\mathbf{\Theta}^{(2)}}^{T}\\
      &=\begin{bmatrix}
        0.5671 & 0.5720 & 0.4428 \\
        0.5548 & 0.5150 & 0.4354
      \end{bmatrix}_{2\times3}\cdot
      \begin{bmatrix}
        0.7  \\
        -0.1 \\
        0.2
      \end{bmatrix}_{3\times1}\\
      &=
      \begin{bmatrix}
        0.4283 \\
        0.4239
      \end{bmatrix}_{2\times1}\\
      \mathbf{A}^{(2)}&= \hat{\mathbf{y}}= \sigma(\mathbf{Z}^{(2)})\\
      &= \sigma\left(\begin{bmatrix}
        0.4283 \\
        0.4239
      \end{bmatrix}_{2\times1}\right)\\
      &=\begin{bmatrix}
        0.6055 \\
        0.6044
      \end{bmatrix}_{2\times1}
        \end{align*}                
        \end{column}
        \begin{column}{.4\textwidth}
            \begin{figure}
                \centering
                \resizebox{\linewidth}{!}{
                \begin{tikzpicture}[remember picture]
    \pgfmathsetmacro{\ex}{1} \pgfmathsetmacro{\ey}{.9} \pgfmathsetmacro{\eradius}{0.9}
    
    \fill (-1, 1) circle[x radius=0.15, y radius=0.15];
    \node at (-1.5, 1) {$x_{1}$};

    \draw[thick, ->] (-1, 1) -- (\ex+1., 0) node[pos=0.6, above]  {$\theta^{(1)}_{21}$};
    \draw[thick, ->] (-1, 1) -- (\ex+1., 2.5) node[pos=0.6, above]  {$\theta^{(1)}_{11}$};
    \draw[thick, ->] (-1, 1) -- (\ex+1., -2.5) node[pos=0.6, above]  {$\theta^{(1)}_{31}$};
    
    \fill (-1, -1) circle[x radius=0.15, y radius=0.15];
    \node at (-1.5, -1) {$x_{2}$};
    \draw[thick, ->] (-1, -1) -- (\ex+1., 2.5) node[pos=0.6, above] {$\theta^{(1)}_{12}$};
    \draw[thick, ->] (-1, -1) -- (\ex+1., 0) node[pos=0.6,above] {$\theta^{(1)}_{22}$};
    \draw[thick, ->] (-1, -1) -- (\ex+1., -2.5) node[pos=0.6, above] {$\theta^{(1)}_{32}$};

    \draw[thick, ->] (\ex+1., 2) -- (\ex+4.5, 0) node[pos=0.41, above]  {$\theta^{(2)}_{11}$};
    \draw[thick, ->] (\ex+1., 0) -- (\ex+6, 0) node[pos=0.3, above]  {$\theta^{(2)}_{12}$};
    \draw[thick, ->] (\ex+1., -2) -- (\ex+4.5, 0) node[pos=0.41, above]  {$\theta^{(2)}_{13}$};

    \draw[fill=PastelBlue] (\ex+1,\ey+2) arc (90:270:\eradius) -- cycle;
    \draw[fill=PastelBlue2] (\ex+1,\ey+2) arc (90:-90:\eradius) -- cycle;
    \node[font=\Large] (z1_1) at (\ex + .6, 0.1+2) {$z^{(1)}_1$};
    \node[font=\Large] (a1_1) at (\ex + 1.5, 0.1+2.) {$a^{(1)}_1$};

    \draw[fill=PastelBlue] (\ex+1,\ey) arc (90:270:\eradius) -- cycle;
    \draw[fill=PastelBlue2] (\ex+1,\ey) arc (90:-90:\eradius) -- cycle;
    \node[font=\Large] (z1_2) at (\ex + .6, 0.1) {$z^{(1)}_2$};
    \node[font=\Large] (a1_2) at (\ex + 1.5, 0.1) {$a^{(1)}_2$};

    \draw[fill=PastelBlue] (\ex+1,\ey-2) arc (90:270:\eradius) -- cycle;
    \draw[fill=PastelBlue2] (\ex+1,\ey-2) arc (90:-90:\eradius) -- cycle;
    \node[font=\Large] (z1_3) at (\ex + .6, 0.1-2) {$z^{(1)}_3$};
    \node[font=\Large] (a1_3) at (\ex + 1.5, 0.1-2.) {$a^{(1)}_3$};

    \draw[fill=PastelBlue] (\ex+4,\ey) arc (90:270:\eradius) -- cycle;
    \draw[fill=PastelBlue2] (\ex+4,\ey) arc (90:-90:\eradius) -- cycle;
    \node[font=\Large] (z2) at (\ex + 3.6, 0.1) {$z^{(2)}_1$};
    \node[font=\Large] (a2) at (\ex + 4.5, 0.1) {$a^{(2)}_1$};
\end{tikzpicture}}
            \end{figure}
        \end{column}    
    \end{columns}
\end{frame}

\begin{frame}{Exemplo Vetorizado: Cálculo da \textit{Loss}}
    \begin{columns}
        \begin{column}{.6\linewidth}
        \small
        \begin{align*}
         J &= \frac{1}{2}\sum (\mathbf{y} - \hat{\mathbf{y}})^{2}\\
         &= \frac{1}{2}\sum\left(\begin{bmatrix}
        0.6055 \\
        0.6044
      \end{bmatrix}_{2\times1} -              
      \begin{bmatrix}
                0.6 \\
                0.8
              \end{bmatrix}_{2\times1}\right)^2\\
      &= 0.01914
        \end{align*}                
        \end{column}
        \begin{column}{.4\textwidth}
            \begin{figure}
                \centering
                \resizebox{\linewidth}{!}{
                \begin{tikzpicture}[remember picture]
    \pgfmathsetmacro{\ex}{1} \pgfmathsetmacro{\ey}{.9} \pgfmathsetmacro{\eradius}{0.9}
    
    \fill (-1, 1) circle[x radius=0.15, y radius=0.15];
    \node at (-1.5, 1) {$x_{1}$};

    \draw[thick, ->] (-1, 1) -- (\ex+1., 0) node[pos=0.6, above]  {$\theta^{(1)}_{21}$};
    \draw[thick, ->] (-1, 1) -- (\ex+1., 2.5) node[pos=0.6, above]  {$\theta^{(1)}_{11}$};
    \draw[thick, ->] (-1, 1) -- (\ex+1., -2.5) node[pos=0.6, above]  {$\theta^{(1)}_{31}$};
    
    \fill (-1, -1) circle[x radius=0.15, y radius=0.15];
    \node at (-1.5, -1) {$x_{2}$};
    \draw[thick, ->] (-1, -1) -- (\ex+1., 2.5) node[pos=0.6, above] {$\theta^{(1)}_{12}$};
    \draw[thick, ->] (-1, -1) -- (\ex+1., 0) node[pos=0.6,above] {$\theta^{(1)}_{22}$};
    \draw[thick, ->] (-1, -1) -- (\ex+1., -2.5) node[pos=0.6, above] {$\theta^{(1)}_{32}$};

    \draw[thick, ->] (\ex+1., 2) -- (\ex+4.5, 0) node[pos=0.41, above]  {$\theta^{(2)}_{11}$};
    \draw[thick, ->] (\ex+1., 0) -- (\ex+6, 0) node[pos=0.3, above]  {$\theta^{(2)}_{12}$};
    \draw[thick, ->] (\ex+1., -2) -- (\ex+4.5, 0) node[pos=0.41, above]  {$\theta^{(2)}_{13}$};

    \draw[fill=PastelBlue] (\ex+1,\ey+2) arc (90:270:\eradius) -- cycle;
    \draw[fill=PastelBlue2] (\ex+1,\ey+2) arc (90:-90:\eradius) -- cycle;
    \node[font=\Large] (z1_1) at (\ex + .6, 0.1+2) {$z^{(1)}_1$};
    \node[font=\Large] (a1_1) at (\ex + 1.5, 0.1+2.) {$a^{(1)}_1$};

    \draw[fill=PastelBlue] (\ex+1,\ey) arc (90:270:\eradius) -- cycle;
    \draw[fill=PastelBlue2] (\ex+1,\ey) arc (90:-90:\eradius) -- cycle;
    \node[font=\Large] (z1_2) at (\ex + .6, 0.1) {$z^{(1)}_2$};
    \node[font=\Large] (a1_2) at (\ex + 1.5, 0.1) {$a^{(1)}_2$};

    \draw[fill=PastelBlue] (\ex+1,\ey-2) arc (90:270:\eradius) -- cycle;
    \draw[fill=PastelBlue2] (\ex+1,\ey-2) arc (90:-90:\eradius) -- cycle;
    \node[font=\Large] (z1_3) at (\ex + .6, 0.1-2) {$z^{(1)}_3$};
    \node[font=\Large] (a1_3) at (\ex + 1.5, 0.1-2.) {$a^{(1)}_3$};

    \draw[fill=PastelBlue] (\ex+4,\ey) arc (90:270:\eradius) -- cycle;
    \draw[fill=PastelBlue2] (\ex+4,\ey) arc (90:-90:\eradius) -- cycle;
    \node[font=\Large] (z2) at (\ex + 3.6, 0.1) {$z^{(2)}_1$};
    \node[font=\Large] (a2) at (\ex + 4.5, 0.1) {$a^{(2)}_1$};
\end{tikzpicture}}
            \end{figure}
        \end{column}    
    \end{columns}

\end{frame}

\begin{frame}{Exemplo Vetorizado: \textit{Backward Pass}}
        \small
        \begin{align*}
      \bm{\delta}^{(2)} &= \left(
      \begin{bmatrix}
        0.6055 \\
        0.6044
      \end{bmatrix}
      -
      \begin{bmatrix}
        0.6 \\
        0.8
      \end{bmatrix}
      \right) \odot
      \begin{bmatrix}
        0.2389 \\
        0.2391
      \end{bmatrix}\\
      &=
      \begin{bmatrix}
        0.0055  \\
        -0.1956
      \end{bmatrix}
      \odot
      \begin{bmatrix}
        0.2389 \\
        0.2391
      \end{bmatrix}\\
      &=
      \begin{bmatrix}
        0.00131 \\
        -0.04676
      \end{bmatrix}\\
      \bm{\delta}^{(1)} &= (\bm{\delta}^{(2)}\cdot \mathbf{\Theta}^{(2)}) \odot \sigma'(\mathbf{Z}^{(1)})\\
      &= \left(
      \begin{bmatrix}
        0.00131 \\
        -0.04676
      \end{bmatrix}
      \cdot
      \begin{bmatrix}
        0.7 & -0.1 & 0.2
      \end{bmatrix}
      \right) \odot \sigma'\left(\begin{bmatrix}
        0.27 & 0.29 & -0.23 \\
        0.22 & 0.06 & -0.26
      \end{bmatrix}\right)\\
      &=
      \begin{bmatrix}
        0.00022 & -0.00003 & 0.00006 \\
        -0.00809 & 0.00117 & -0.00230
      \end{bmatrix}
        \end{align*}                
\end{frame}

 \begin{frame}{Exemplo Vetorizado: \textit{Backward Pass}}
    \begin{columns}
        \begin{column}{.6\linewidth}
        \small
        \begin{align*}
              \frac{\partial J}{\partial \mathbf{\Theta}^{(2)}}&={\bm{\delta}^{(2)}}^{T}\cdot \mathbf{A}^{(1)}\\
              &=
      \begin{bmatrix}
        0.00131 & -0.04676
      \end{bmatrix}
      \cdot
      \begin{bmatrix}
        0.5671 & 0.5720 & 0.4428 \\
        0.5548 & 0.5150 & 0.4354
      \end{bmatrix}\\
            &=
      \begin{bmatrix}
        -0.0252 & -0.0233 & -0.0198
      \end{bmatrix}\\
      \frac{\partial J}{\partial \mathbf{\Theta}^{(1)}}&={\bm{\delta}^{(1)}}^{T}\cdot \mathbf{X} \\
            &=
      \begin{bmatrix}
        0.00022 & -0.00809 \\
        -0.00003 & 0.00117 \\
        0.00006 & -0.00230 \\
      \end{bmatrix}\cdot \begin{bmatrix}
                0.5 & 0.1 \\
                0.2 & 0.6
              \end{bmatrix}\\
      &=
      \begin{bmatrix}
        -0.00150 & -0.00483 \\
        0.00022  & 0.00070  \\
        -0.00043 & -0.00137
      \end{bmatrix}
        \end{align*}                
        \end{column}
        \begin{column}{.4\textwidth}
            \begin{figure}
                \centering
                \resizebox{\linewidth}{!}{
                \begin{tikzpicture}[remember picture]
    \pgfmathsetmacro{\ex}{1} \pgfmathsetmacro{\ey}{.9} \pgfmathsetmacro{\eradius}{0.9}
    
    \fill (-1, 1) circle[x radius=0.15, y radius=0.15];
    \node at (-1.5, 1) {$x_{1}$};

    \draw[thick, ->] (-1, 1) -- (\ex+1., 0) node[pos=0.6, above]  {$\theta^{(1)}_{21}$};
    \draw[thick, ->] (-1, 1) -- (\ex+1., 2.5) node[pos=0.6, above]  {$\theta^{(1)}_{11}$};
    \draw[thick, ->] (-1, 1) -- (\ex+1., -2.5) node[pos=0.6, above]  {$\theta^{(1)}_{31}$};
    
    \fill (-1, -1) circle[x radius=0.15, y radius=0.15];
    \node at (-1.5, -1) {$x_{2}$};
    \draw[thick, ->] (-1, -1) -- (\ex+1., 2.5) node[pos=0.6, above] {$\theta^{(1)}_{12}$};
    \draw[thick, ->] (-1, -1) -- (\ex+1., 0) node[pos=0.6,above] {$\theta^{(1)}_{22}$};
    \draw[thick, ->] (-1, -1) -- (\ex+1., -2.5) node[pos=0.6, above] {$\theta^{(1)}_{32}$};

    \draw[thick, ->] (\ex+1., 2) -- (\ex+4.5, 0) node[pos=0.41, above]  {$\theta^{(2)}_{11}$};
    \draw[thick, ->] (\ex+1., 0) -- (\ex+6, 0) node[pos=0.3, above]  {$\theta^{(2)}_{12}$};
    \draw[thick, ->] (\ex+1., -2) -- (\ex+4.5, 0) node[pos=0.41, above]  {$\theta^{(2)}_{13}$};

    \draw[fill=PastelBlue] (\ex+1,\ey+2) arc (90:270:\eradius) -- cycle;
    \draw[fill=PastelBlue2] (\ex+1,\ey+2) arc (90:-90:\eradius) -- cycle;
    \node[font=\Large] (z1_1) at (\ex + .6, 0.1+2) {$z^{(1)}_1$};
    \node[font=\Large] (a1_1) at (\ex + 1.5, 0.1+2.) {$a^{(1)}_1$};

    \draw[fill=PastelBlue] (\ex+1,\ey) arc (90:270:\eradius) -- cycle;
    \draw[fill=PastelBlue2] (\ex+1,\ey) arc (90:-90:\eradius) -- cycle;
    \node[font=\Large] (z1_2) at (\ex + .6, 0.1) {$z^{(1)}_2$};
    \node[font=\Large] (a1_2) at (\ex + 1.5, 0.1) {$a^{(1)}_2$};

    \draw[fill=PastelBlue] (\ex+1,\ey-2) arc (90:270:\eradius) -- cycle;
    \draw[fill=PastelBlue2] (\ex+1,\ey-2) arc (90:-90:\eradius) -- cycle;
    \node[font=\Large] (z1_3) at (\ex + .6, 0.1-2) {$z^{(1)}_3$};
    \node[font=\Large] (a1_3) at (\ex + 1.5, 0.1-2.) {$a^{(1)}_3$};

    \draw[fill=PastelBlue] (\ex+4,\ey) arc (90:270:\eradius) -- cycle;
    \draw[fill=PastelBlue2] (\ex+4,\ey) arc (90:-90:\eradius) -- cycle;
    \node[font=\Large] (z2) at (\ex + 3.6, 0.1) {$z^{(2)}_1$};
    \node[font=\Large] (a2) at (\ex + 4.5, 0.1) {$a^{(2)}_1$};
\end{tikzpicture}}
            \end{figure}
        \end{column}    
    \end{columns}

\end{frame}


  \begin{frame}
    \frametitle{Exemplo Vetorizado: \textit{Optimizer}}
    \small
    \textbf{Atualização dos Pesos ($\eta = 0.1$):}
    \[
      \mathbf{\Theta}_{t}^{(l)}=\mathbf{\Theta}_{t-1}^{(l)}-\eta\nabla_{\thetav}J
    \]
    \vfill
    \textbf{Novos Pesos da Camada 2:}
    \begin{align*}
      \mathbf{\Theta}_{new}^{(2)}&=
      \begin{bmatrix}
        0.7 & -0.1 & 0.2
      \end{bmatrix}
      - 0.1 \cdot
      \begin{bmatrix}
        -0.0252 & -0.0233 & -0.0198
      \end{bmatrix}\\
      &=
      \begin{bmatrix}
        0.7025 & -0.0977 & 0.2020
      \end{bmatrix}
    \end{align*}
    \vfill
    \textbf{Novos Pesos da Camada 1:}
    \begin{align*}
      \mathbf{\Theta}_{new}^{(1)}&=
      \begin{bmatrix}
        0.5  & 0.2  \\
        0.6  & -0.1 \\
        -0.4 & -0.3
      \end{bmatrix}
      - 0.1 \cdot
           \begin{bmatrix}
        -0.00150 & -0.00483 \\
        0.00022  & 0.00070  \\
        -0.00043 & -0.00137
      \end{bmatrix}\\
      &=
      \begin{bmatrix}
        0.5002  & 0.2005   \\
        0.6000  & -0.1001 \\
        -0.4000 & -0.2999
      \end{bmatrix}
    \end{align*}
  \end{frame}

  \begin{frame}
    \frametitle{Resumindo}
    \begin{itemize}
      \item Podemos treinar redes neurais de múltiplas camadas usando descida de
        gradiente.
        \vfill

      \item O processo de treinamento consiste em repetir 4 etapas:
        \begin{enumerate}
          \item Computar todas as saídas da rede (\textit{Forward Pass})

          \item Computar o quão diferente as saídas são em relação a variável alvo (\textit{Loss Function})

          \item Computar o gradiente do erro em relação aos parâmetros da rede (\textit{Backward Pass})

          \item Atualizar os pesos (\textit{Optimizer})
        \end{enumerate}
        \vfill

      \item Todo o treinamento pode ser implementado de forma vetorizada.
    \end{itemize}
  \end{frame}

  \section*{Referências e Leituras}
  %------------------------------------------------
  \begin{frame}{Leituras Recomendadas}
    \leiturasize
    \begin{itemize}
      \item Primeira publicação do Backpropagation: \fullcite{rumelhart1986learningbackpropextended}

      \item Capítulo 2: \fullcite{nielsen2015neural}
    \end{itemize}
  \end{frame}

  %------------------------------------------------

  \begin{frame}{Referências}
    \printbibliography
  \end{frame}
\end{document}